% Lliçó: funcions contínues (valencià).

\begin{frame}
\didactatitle{Funcions contínues}

\begin{definition}[Continuïtat]
Siga $f$ una funció definida en un interval $I$ i siga $a \in I$. Diem que
$f$ és \keyterm{contínua en $a$} si per a cada $\varepsilon > 0$ existeix un
$\delta > 0$ tal que
\[
  x \in I, \ |x - a| < \delta \implies |f(x) - f(a)| < \varepsilon .
\]
És \keyterm{contínua en $I$} si ho és en tots els seus punts.
\end{definition}

Si $a$ no és un extrem de $I$, la condició diu exactament que
\begin{keyformula}
  \lim_{x \to a} f(x) = f(a) ,
\end{keyformula}
\bymedium{%
  és a dir: el límit existeix i val el mateix que la funció.
}{%
  és a dir: el límit existeix, $f$ està definida en $a$, i els dos valors
  coincideixen. Si falla qualsevol de les tres coses, $f$ no és contínua en
  $a$.
}
\end{frame}

\begin{frame}
\didactatitle{Operacions i discontinuïtats}

\begin{remark}
Per l'àlgebra de límits, la suma, el producte i el quocient (on el
denominador no s'anul·la) de funcions contínues són continus. En particular,
els polinomis i les funcions racionals són continus allà on estan definits.
\end{remark}

\dpause

\begin{example}
\begin{itemize}
  \item $f(x) = \dfrac{x^2 - 1}{x - 1}$ per a $x \neq 1$, amb $f(1) = 0$. El
        límit en $1$ existeix i val $2$, però no coincideix amb $f(1)$. La
        discontinuïtat és \keyterm{evitable}: n'hi ha prou de redefinir
        $f(1) = 2$.
  \item $H(x) = 0$ per a $x < 0$ i $H(x) = 1$ per a $x \geq 0$. No és
        contínua en $0$: amb $\varepsilon = 1/2$, qualsevol interval al
        voltant de $0$ conté punts $x < 0$, on $|H(x) - H(0)| = 1$. És un
        \keyterm{salt}, i cap valor de $H(0)$ no l'arregla.
\end{itemize}
\end{example}

\begin{commonmistake}
Confonen «contínua» amb «es dibuixa sense alçar el llapis», i no saben què
fer amb $1/x$, que és contínua en $(-\infty, 0)$ i en $(0, +\infty)$ encara
que la seua gràfica tinga dues branques.
\end{commonmistake}
\end{frame}
